\section{Actual motion, ramified branches, and contour traces}

\begin{theorem}[NN-A004: all shifted-series jets and factorization obstruction]
Let $L(s)=\sum_{n\ge1}a_n n^{-s}$ converge absolutely for $\Re s>1$
and continue holomorphically near $\rho$, where $\Re\rho>0$ and
$L(\rho)=0$. In a sufficiently small neighborhood of $(\rho,0)$, set
\[
 L_t(s)=\sum_{k\ge0}\frac{(-1)^k(s)_k}{k!}L(s+k)t^k.
\]
This is holomorphic and agrees with $\sum a_n(n+t)^{-s}$ in the initial
half-plane. If $\rho$ is simple, its unique analytic zero branch has
\[
 v=\rho'(0)=\frac{\rho L(\rho+1)}{L'(\rho)},
\quad
 \rho''(0)=-\frac{L''(\rho)v^2
 -2\{L(\rho+1)+\rho L'(\rho+1)\}v
 +\rho(\rho+1)L(\rho+2)}{L'(\rho)}.
\]
All higher coefficients $\rho(t)=\rho+\sum_{n\ge1}c_nt^n$ are determined
without a new existence hypothesis by
\[
 c_n=-\frac{[t^n]L_t(\rho+\sum_{j<n}c_jt^j)}{L'(\rho)}.
\]
For real $t\downarrow0$, the first nonzero $\Re c_n$ determines the
normal side of the critical line. If every $\Re c_n=0$, the branch stays
on that line near zero. A nonzero complex velocity alone does not exclude
normal tangency.

If $L=FG$, $F(\rho)=0$, $F'(\rho)G(\rho)\ne0$, then the shifts of the
Dirichlet coefficients satisfy
\[
 v_L=v_F\,\frac{G(\rho+1)}{G(\rho)}.
\]
Consequently factorwise GRH does not give a factorwise rule for the motion
labels: the nonvanishing factor has a retained complex multiplier.
\end{theorem}
\begin{proof}
For $s$ in a fixed compact neighborhood of $\rho$, choose $k_0$ such that
$\Re(s+k_0)>1$. Absolute convergence gives a uniform bound for $L(s+k)$
for $k\ge k_0$, while $(s)_k/k!=O(k^C)$ uniformly on this compact set.
The displayed series and its local holomorphic derivatives converge for
$|t|<1$. There are only finitely many lower $k$, and their neighborhoods
can avoid any poles. Expanding $(n+t)^{-s}$ by the binomial theorem and
using absolute convergence proves equality in the initial half-plane.
The implicit function theorem applies because $L'(\rho)\ne0$.
The first two derivatives of $L_t(\rho(t))=0$ use
$\partial_tL_t|_0=-sL(s+1)$ and
$\partial_t^2L_t|_0=s(s+1)L(s+2)$, and give the stated formulas.
Adding $c_nt^n$ to the already constructed truncation changes its degree
$n$ coefficient by exactly $L'(\rho)c_n$, proving the recursion. A
real analytic function is either identically zero locally or has the sign
of its first nonzero coefficient for small positive argument. Finally
$L'(\rho)=F'(\rho)G(\rho)$ and
$L(\rho+1)=F(\rho+1)G(\rho+1)$ prove the product formula.
\end{proof}

This applies to $\zeta$, Dirichlet $L$-functions, and Dedekind $\zeta_K$.
A supplies the critical line for $\zeta$ and Dirichlet factors; for abelian
$K/\mathbb Q$ the standard Dirichlet factorization transfers that location
to $\zeta_K$. General $K$ requires the stipulated Hecke input M, or a
separate GRH hypothesis. None of these inputs supplies simplicity. At a
nontrivial critical zero, the Euler product at $\rho+1$ shows $v\ne0$
when the zero is simple; it gives no sign for $\Re v$.
The geometric object here is the analytic zero divisor in the two
coordinates $(s,t)$ and its sections over the parameter disc. It is not
an algebraic family over a finite field, and its complex multipliers have
not been identified with Frobenius eigenvalues.

\begin{theorem}[NN-A005: multiple-zero branches and the residue trace]
Suppose instead that $L$ has multiplicity $m\ge2$ at $\rho$ and
$\rho L(\rho+1)\ne0$. Put
\[
 C=\frac{m!\rho L(\rho+1)}{L^{(m)}(\rho)}.
\]
The $m$ local branches on the ramified parameter $u=t^{1/m}$ satisfy
\[
 s_j(u)=\rho+c\omega_m^j u+O(u^2),\qquad c^m=C.
\]
In particular they are not individually differentiable as holomorphic
functions of $t$ at zero. Their symmetric trace is analytic. More
generally, let $\Gamma$ be a positively oriented finite union of simple
contours, bounding a region where $L_t$ is holomorphic and has no boundary
zeros for small $t$. If $P$ is holomorphic on that region, then
\[
 \left.\frac d{dt}\sum_{L_t(z)=0}P(z)\right|_{t=0}
 =\frac1{2\pi i}\int_\Gamma
 P'(s)\frac{sL(s+1)}{L(s)}\,ds.
\]
Zeros in this formula are counted with multiplicity. For $P(s)=s$, the
contribution of a multiplicity-$m$ zero is
\[
 \operatorname{Res}_{s=\rho}\frac{sL(s+1)}{L(s)}
 =\frac1{(m-1)!}
 \left.\left(\frac d{ds}\right)^{m-1}
 \frac{sL(s+1)}{L(s)/(s-\rho)^m}\right|_{s=\rho}.
\]
Thus a contour mean remains meaningful without simple zeros, whereas a
sum of the simple-zero velocities does not.
\end{theorem}
\begin{proof}
With $a=L^{(m)}(\rho)/m!$, the quotient
$u^{-m}L_{u^m}(\rho+uy)$ extends holomorphically to $u=0$ and there
equals $ay^m-\rho L(\rho+1)$. Its $m$ roots are nonzero and simple.
The implicit function theorem gives the $m$ analytic $y_j(u)$, and
Rouche's theorem counts all roots in the local cluster. The nonzero
leading terms $c\omega_m^j t^{1/m}$ exclude holomorphic individual
branches in $t$.
The argument principle writes the symmetric sum as
$(2\pi i)^{-1}\int_\Gamma P(s)\partial_s\log L_t(s)\,ds$.
The quotient $\partial_tL_t/L_t$ is analytic in $t$ on the compact
contour, so differentiation under the integral and integration by parts
give $-(2\pi i)^{-1}\int P'(s)\partial_tL_t/L_t$. Insert
$\partial_tL_t|_0=-sL(s+1)$. The residue expression is the ordinary
formula for a pole of order at most $m$.
\end{proof}

For $L=\zeta$, the right-hand contour integrand has the fully explicit
Dirichlet series, initially on $\Re s>1$,
\[
 \frac{s\zeta(s+1)}{\zeta(s)}
 =s\sum_{n\ge1}\frac{b(n)}{n^s},\qquad
 b(n)=\frac1n\prod_{p\mid n}(1-p).
\]
Indeed convolution with $1/\zeta(s)=\sum\mu(n)n^{-s}$ gives
$b(n)=n^{-1}\sum_{d\mid n}d\mu(d)$. For general $L$, replace
$b(n)$ by $\sum_{d\mid n}a_d d^{-1}a^{-1}_{n/d}$ whenever the
Dirichlet inverse converges absolutely. This is an exact residue
formula, not an unproved asymptotic on contours passing through zeros.
In the zero-divisor geometry, symmetric sums are functions on the relative
degree-$m$ divisor (or its symmetric power); individual sections require
the stated ramified base change. No complexity bound follows from the
existence of the contour integral.

\begin{proposition}[NN-A010: the completed Hurwitz deformation has extra poles]
Let $\Lambda_t(s)=\pi^{-s/2}\Gamma(s/2)\zeta(s,1+t)$. For $m\ge1$,
\[
 \operatorname{Res}_{s=-2m}\Lambda_t(s)
 =-\frac{2(-1)^m\pi^m}{m!(2m+1)}B_{2m+1}(1+t).
\]
At $t=0$ these residues vanish; their first $t$ derivatives are
$-2(-1)^m\pi^mB_{2m}/m!\ne0$. At $t=1$, every such residue is
$-2(-1)^m\pi^m/m!$. Consequently the rational principal-part
bookkeeping using the poles at zero and one is not the complete
meromorphic divisor of $\Lambda_t$. It does not provide a meromorphic
function on $\mathbb P^1$ for the full completion.
\end{proposition}
\begin{proof}
The residue of $\Gamma(s/2)$ at $s=-2m$ is $2(-1)^m/m!$.
Use the exact Hurwitz value
$\zeta(-2m,a)=-B_{2m+1}(a)/(2m+1)$ (DLMF 25.11.14).
Bernoulli differentiation and $B_{2m}(1)=B_{2m}$ give the derivative.
For $t=1$, $\zeta(s,2)=\zeta(s)-1$, so its value at every negative
even integer is $-1$. A meromorphic function on compact
$\mathbb P^1$ has finitely many poles, whereas this example has
infinitely many. This proves the asserted scope restriction.
\end{proof}
